%-------------------------------------------------------------------------------
% IPOL LaTeX class test example
% by rafael grompone von gioi
% June 4, 2012
%-------------------------------------------------------------------------------
\documentclass{ipol}

\ipolSetTitle{IPOL \LaTeX\ AD Census}
\ipolSetAuthors{Rafael Grompone von Gioi$^1$ \& Rapha\"el Toucour$^2$}
\ipolSetAffiliations{%
$^1$ CMLA, ENS Cachan, France (\texttt{grompone@cmla.ens-cachan.fr})\\
$^2$ IIE, UdelaR, Uruguay (\texttt{jirafa@fing.edu.uy})}

\begin{document}

%-------------------------------------------------------------------------------
\begin{abstract}
This document describes and illustrates how to use the IPOL's
\LaTeX\ class, created to produces a uniform layout for IPOL's
articles. The restrictions imposed on articles to be published in IPOL
are briefly discussed.
\end{abstract}

%-------------------------------------------------------------------------------
\section{Getting Started}

There is no need for installation to use this class, just two files
need to be copied to the same directory were are the \LaTeX\ source
files. These files are the class itself, \verb|ipol.cls|, and the logo
file, that can be either \verb|ipol_logo.eps| if you compile with
\verb|latex| or \verb|ipol_logo.png| if you compile with
\verb|pdflatex|. If not sure, just copy the three of these files.

The minimal example of the class is as follows:
\begin{verbatim}
\documentclass{ipol}
\begin{document}
\end{document}
\end{verbatim}
It will only generates IPOL's header, including it's logo, and the
words ``title'' and ``authors'' where the title and authors should be
placed.  This example is useless but can be used to test your system.

This class is based on the standard 'article' class of \LaTeX\ and it
is used essentially in the same way. There are two main use
instruction: the layout must not be changed and title is not generated
with the usual \verb|\title|, \verb|\author|, \verb|\date|, and
\verb|\maketitle| commands. Let us discuss one by one.

This class was created to provide a uniform layout for IPOL articles,
so no command should be used that would change the page layout: do not
change paper size (it must be A4 paper), do not change the page
margins, do not change the font type, size, or color.

The commands \verb|\ipolSetTitle|, \verb|\ipolSetAuthors| and
\verb|\ipolSetAffiliations|, are provided by the IPOL class to set the
information needed to generate the title of the article. These
commands must be used in the preamble of the \LaTeX\ file, that is
between the \verb|\documentclass{ipol}| and \verb|\begin{document}|
commands. This is important, otherwise the title will not be generated
correctly. As the name of the commands imply, \verb|\ipolSetTitle| is
used to set the title of the article, \verb|\ipolSetAuthors| is used
to set the article's authors, and \verb|\ipolSetAffiliations| is used
to set the authors affiliations. This last command is optional. The
following is an example of how to use it:
\begin{verbatim}
\documentclass{ipol}
\ipolSetTitle{Ceci n'est pas un article}
\ipolSetAuthors{Rafael Grompone von Gioi}
\ipolSetAffiliations{CMLA, ENS Cachan, France}
\begin{document}
\end{document}
\end{verbatim}
If there are more than one author, and they have different affiliations,
they may be indicated as follows:
\begin{verbatim}
\documentclass{ipol}
\ipolSetTitle{Ceci n'est pas un article}
\ipolSetAuthors{Rafael Grompone von Gioi$^1$ \& Raphael Toucour$^2$}
\ipolSetAffiliations{$^1$ CMLA, ENS Cachan, France\\
                     $^2$ IIE, UdelaR, Uruguay}
\begin{document}
\end{document}
\end{verbatim}
For the rest, the class is used as the standard ``article'' class:
packages or definitions can be added in the preamble, and the main
text can contain sections, subsections, mathematics, figures, tables,
etc.

The rest of this document describes some special command for IPOL and
briefly comments on some conditions that will be imposed for a final
article, just before publishing at IPOL. (These restrictions are not
necessary for submissions but will reduce the work later.)

%-------------------------------------------------------------------------------
\section{Specific IPOL Commands}

Once an article is published at IPOL it will be assigned a DOI number
and a universal DOI web address to refer to it. The IPOL class
provides a command that will later contain the article's DOI number
(\verb|\ipolDOI|) and a command that will contain the article's web
address (\verb|\ipolLink|). These are the only reference that should
be used for the IPOL article itself and to its complementary contents
on the web.

These commands will have a default value in a preprint version, but
once the paper is accepted they will be replaced by the right value.
One may want to have a valid link while the article is being prepared
and peer reviewed. For that purpose, the command \verb|\ipolPreprintLink|
is used to set the article link \verb|\ipolLink| to a temporary
value. For example,
\begin{verbatim}
\ipolPreprintLink{https://tools.ipol.im/wiki/ref/manuscript_guidelines/}
\end{verbatim}
This command should be used in the document's preamble (i.e., after
\verb|\documentclass{ipol}| and before \verb|\begin{document}|).

The command \verb|\ipolLink| only provides the address of the
corresponding IPOL article. To produce a link with it or other
address, the command \verb|\href{}{}| should be used. The first
parameter is the link address while the second is the text that will
be shown in the resulting document. Also, to be able to read the
address in a printed version of the document, the address is also
printed as footnote. For example, the commands
\begin{verbatim}
\href{\ipolLink}{the article page}
\href{https://tools.ipol.im/wiki/ref/manuscript_guidelines/}{guidelines}
\end{verbatim}
produce the following links:
\href{\ipolLink}{the article page} and
\href{https://tools.ipol.im/wiki/ref/manuscript_guidelines/}{guidelines}.

Two particular environments are defined to generate two particular
parts of IPOL articles: \verb|ipolCode| and \verb|ipolSupp|. The first
one generates the ``Source Code'' description paragraph while the
second generates the ``Supplementary Material'' description
paragraph. These environment work similarly as the ``abstract'' and
should be used after the \verb|\begin{document}| and before the main
text of the article. The following is an example:
\begin{verbatim}
\documentclass{ipol}
\ipolSetTitle{Ceci n'est pas un article}
\ipolSetAuthors{Rafael Grompone von Gioi}
\begin{document}

\begin{abstract}
A short description of the article.
\end{abstract}

\begin{ipolCode}
Description of the source code to be found \href{\ipolLink}{here}.
\end{ipolCode}

\begin{ipolSupp}
Some articles could provide additional material, not part of the
peer reviewed article but related and useful. It should be found at
\href{\ipolLink}{the web page}.
\end{ipolSupp}

\section{Introduction}
The article main text starts here.

\end{document}
\end{verbatim}

%-------------------------------------------------------------------------------
\section{IPOL final article restrictions}

This document describes IPOL \LaTeX\ class that provides a uniform
layout for IPOL articles. Authors should also check the
\href{https://tools.ipol.im/wiki/ref/manuscript_guidelines/}{IPOL
  manuscript guidelines} where the requirements and suggestions about
the contents, style and description detail for IPOL articles are
described.

The PDF files for IPOL submissions are generated by the authors using
any method that produces a readable PDF file. Once an article is
accepted, however, the final version will be generated at IPOL servers
in order to provide PDF files with controlled sizes and resolutions
conditions. For that aim, there are some restrictions that must be
imposed to the source files. In what follows are described the main
restrictions, but more details are described
\href{https://tools.ipol.im/wiki/ref/manuscript_guidelines/}{here}.

The final documents will be generated using pdf\LaTeX. This forces
some restrictions, in particular the only graphical formats directly
accepted are PNG, JPEG and PDF itself. Images or figures in other kind
of formats should be converted to one of these.

The files and directories must follow a strict naming as defined
\href{https://tools.ipol.im/wiki/ref/manuscript_guidelines/}{here}.

Finally, the authors must warrant to be the authors of every part
of the final article, or must include permission of the copyright
holder. In particular, if images not owned by the authors are used,
the authors and permissions to be used must be included in the article.
For example, by adding a final section as follows:
\begin{verbatim}
\section*{Image Credits}

\includegraphics[height=2em]{image1.png}
\href{http://a.link.here}{Courtesy of Name Surname}
\\
\includegraphics[height=2em]{image2.png}
\copyright\ Pierre Dupont
\\
\includegraphics[height=2em]{image3.png}
the authors
\end{verbatim}

%-------------------------------------------------------------------------------
\section*{Image Credits}

\includegraphics[height=2em]{ipol_logo} \copyright\ IPOL (there's no
need to credit this image, here is used as an example.)

\end{document}
%-------------------------------------------------------------------------------
